\section*{第 \ref{ch:g9:alternator} 章——从电站到家：交流发电机}

\begin{solution}{exo:g9:alternator:1}
只有变化才会产生感应；运动更快、\omterm{def:g1:magnets:magnet}{磁铁}更强、
匝数更多，\omterm{def:g8:voltage-voltmeter:voltage}{电压}就更大；运动反向，电流也随之反向。
静止的\omterm{def:g1:magnets:magnet}{磁铁}，无论多强，
都不会产生感应。
\end{solution}

\begin{solution}{exo:g9:alternator:2}
转子每转半圈，\omterm{def:g4:magnets-and-compass:poles}{磁极}便靠近线圈（感应出一个方向的\omterm{def:g8:voltage-voltmeter:voltage}{电压}），
随后又将其带离（感应出相反方向的\omterm{def:g8:voltage-voltmeter:voltage}{电压}）：
\omterm{def:g8:voltage-voltmeter:voltage}{电压}必须每半圈变换一次正负——
交变性本就内含于旋转之中。
\end{solution}

\begin{solution}{exo:g9:alternator:3}
每秒五十转（对于简单的双极机）。
每个电站都必须保持完全相同的转动节奏——
整个电网同步旋转，共享同一\omterm{def:g8:sound-pitch-loudness:frequency}{频率}。
\end{solution}

\begin{solution}{exo:g9:alternator:4}
涡轮机与\omterm{def:g9:alternator:alternator}{交流发电机}。只有推动叶片的动力
会变化：火焰产生的蒸汽、原子能产生的蒸汽、落下的湖水，或风。
\end{solution}

\begin{solution}{exo:g9:alternator:5}
湖泊的储存高度 $\to$ 落水的运动 $\to$
涡轮机的转动 $\to$ 发电机的电能 $\to$
变压器与电网 $\to$ \omterm{def:g3:first-electric-circuit:bulb}{灯泡}。太阳为湖泊
注入了\omterm{def:g5:everyday-energy:energy}{能量}：\omterm{def:g2:ice-water-steam:evapcond}{蒸发}将海水抬升，
雨水再把它送到山区。
\end{solution}

\begin{solution}{exo:g9:alternator:6}
\omterm{def:g3:first-electric-circuit:bulb}{灯泡}的\omterm{def:g5:everyday-energy:energy}{能量}必须来自某处，而自行车上唯一的
储存\omterm{def:g5:everyday-energy:energy}{能量}的地方就是骑车人：发电机把腿部做的功
转化为光，因此\omterm{def:g3:first-electric-circuit:bulb}{灯泡}的\omterm{def:g9:kinetic-energy-safety:joule}{焦耳}数表现为额外的阻力。
\omterm{def:g5:everyday-energy:energy}{能量}只会转化，绝不会凭空产生——
这是本课程最古老的定律。
\end{solution}

\begin{solution}{exo:g9:alternator:7}
\omterm{ex:g9:alternator:chains}{可再生的}：风、湖泊（雨水会补充它）、阳光。不\omterm{ex:g9:alternator:chains}{可再生的}：
煤、铀——以人类的时间尺度来看，用掉便永远没有了。
分类的依据是：我们消耗储存时，自然界能否将其补充？
\end{solution}

\begin{solution}{exo:g9:alternator:8}
$P = U I$ 表明，同样的\omterm{def:g9:electric-power-energy:power}{功率}可以以高\omterm{def:g8:voltage-voltmeter:voltage}{电压}、小电流的形式输送；
而\omterm{prop:g8:electric-current-effects:three}{热效应}使导线的损耗随电流增大而增加。
电流小，导线便不易发热，\omterm{def:g5:everyday-energy:energy}{能量}得以节省——
因此电网以高压输送，到了用户门口再降压。
\end{solution}

\begin{solution}{exo:g9:alternator:9}
慢速摇动：感应\omterm{def:g8:voltage-voltmeter:voltage}{电压}小（规则二）——灯光微弱，
每秒\omterm{def:g9:alternating-current:period}{周期}数很少：出现肉眼可见的闪烁，
低于人眼的融合速率。快速摇动：\omterm{def:g8:voltage-voltmeter:voltage}{电压}更大——
灯光明亮，\omterm{def:g8:sound-pitch-loudness:frequency}{频率}高到人眼看来平滑连续。两个旋钮，
一根摇柄。
\end{solution}

\begin{solution}{exo:g9:alternator:10}
在 \qty{230}{V} 下：$I = \num{1e8} \div 230 \approx
\qty{4.3e5}{A}$——四十万\omterm{def:g8:intensity-ammeter:intensity}{安培}：地球上没有任何电缆能承受。
在 \qty{400000}{V} 下：$I = \num{1e8} \div
\num{4e5} = \qty{250}{A}$——普通输电电缆即可。
变压器的问题就此得到解答。
\end{solution}

\begin{solution}{exo:g9:alternator:11}
两条链条都是：储存的\omterm{def:g5:everyday-energy:energy}{能量} $\to$ 热 $\to$ 沸腾的水
$\to$ 蒸汽的运动 $\to$ 转动 $\to$ 电能——
精巧的水壶为汽轮机提供蒸汽。两条链条在
恰好一个环节上不同：即产生\omterm{def:g5:heat-and-insulation:heat}{热量}的方式——燃烧煤的化学反应，
或裂变铀的原子核。
\end{solution}

\begin{solution}{exo:g9:alternator:12}
向上追溯：墙上插座的 \qty{230}{V} 经过
社区变压器（\omterm{def:g8:voltage-voltmeter:voltage}{电压}经降压），沿
数百千伏的高压塔线路（输送形式：高压\omterm{def:g9:alternating-current:dcac}{交流电}，
略有\omterm{def:g5:heat-and-insulation:heat}{热量}散失）而来，来自一台\omterm{def:g9:alternator:alternator}{交流发电机}（依靠感应，把转动变为电能），
由汽轮机带动（蒸汽或水的运动），
而汽轮机的动力来自某种储存——比如水坝抬高的湖泊，
湖泊又是由太阳的\omterm{def:g2:ice-water-steam:evapcond}{蒸发}抬升的。所以九年级那句话：
你房间里的灯光就是阳光——
被储存、下落、转动，然后输送而来。
\end{solution}

\begin{solution}{pb:g9:alternator:1}
\textbf{1.} 湖泊的高度 $\to$ 水流奔涌的运动
$\to$ 涡轮机的转动 $\to$ \omterm{def:g9:alternator:alternator}{交流发电机}感应出的
交变\omterm{def:g8:voltage-voltmeter:voltage}{电压} $\to$ 电站的母线。
\textbf{2.} 转子\omterm{def:g1:magnets:magnet}{磁铁}转动，线圈固定不动。
没有运动，穿过线圈的磁场就不会变化——
而没有变化，感应便毫无作用。
\textbf{3.} \omterm{def:g8:voltage-voltmeter:voltage}{电压}的\omterm{def:g8:sound-pitch-loudness:frequency}{频率}与转动同步：
转速不对，\omterm{def:g8:sound-pitch-loudness:frequency}{频率}就不对——而一台与
电网的 \qty{50}{Hz} 不同步的发电机，
不但无助于电网，反而与网上其他发电机相互抗衡。
\textbf{4.} 更强的转子\omterm{def:g1:magnets:magnet}{磁铁}确实能感应出更多电能——
但涡轮机随即必须更费力，才能在更强的
磁性阻力下维持转动：输出的电能越多，
需要通过的水量也越多，\omterm{def:g9:kinetic-energy-safety:joule}{焦耳}对\omterm{def:g9:kinetic-energy-safety:joule}{焦耳}。
付出代价的是湖泊，而不是\omterm{def:g1:magnets:magnet}{磁铁}。
\textbf{5.} $I = \num{2e7} \div \num{2e5} = \qty{100}{A}$.
\textbf{6.} $I = \num{2e7} \div 230 \approx \qty{8.7e4}{A}$
——八万七千\omterm{def:g8:intensity-ammeter:intensity}{安培}：不可能有这样的电缆。
变压器的作用正是这种交换——\omterm{def:g8:voltage-voltmeter:voltage}{电压}升高，
电流降低，\omterm{def:g9:electric-power-energy:power}{功率}不变。
\textbf{7.} 闸门通过的水量需要增加。水量不足时，
负载下的机组会像上坡的骑车人一样减速——而
电网\omterm{def:g8:sound-pitch-loudness:frequency}{频率}会降到五十\omterm{def:g8:sound-pitch-loudness:frequency}{赫兹}以下：电网上的时钟与
各类机器都会感受到电网的“饥饿”。
\textbf{8.} 只有第一个方框——即风、叶片的
推力。涡轮机、\omterm{def:g9:alternator:alternator}{交流发电机}和变压器都完好无损，
状态良好，只是无风。
\textbf{9.} 电网是一个巨大的\omterm{def:g5:series-parallel-circuits:parallel}{并联}网络：许多
线路将每座电站与每个城镇相连，因此断开的线路的
负荷可以经由未断开的\omterm{def:g5:series-parallel-circuits:parallel}{支路}转移——\omterm{def:g5:series-parallel-circuits:parallel}{并联}的
独立性在全国尺度上得以体现。
\textbf{10.} 改道后、过载的线路使
输送的\omterm{def:g8:voltage-voltmeter:voltage}{电压}一分钟内降到\omterm{def:g8:voltage-voltmeter:nominal}{额定电压}以下：仓库的
灯光变暗，\omterm{ex:g6:circuits-and-safety:motor}{电动机}也转得无精打采——
供电不足，任何设备在低于\omterm{def:g8:voltage-voltmeter:nominal}{额定电压}时都会如此。
\textbf{11.} “今晚的每一瓦电能都来自下落的湖水；
太阳\omterm{def:g2:ice-water-steam:evapcond}{蒸发}了海洋，雨水把它带到
我们的山区。我们在这里转动阳光——阳光储存于
湖泊，经由天空中转。”
\textbf{12.} 例如：“运动经过\omterm{def:g1:magnets:magnet}{磁铁}，造出了每一伏夜间的\omterm{def:g8:voltage-voltmeter:voltage}{电压}——
而到了七点，一万只水壶
都催着湖水再快一点落下。”
\end{solution}
